% PAGES: 77-80
% Line 1 of sec02_c4_1.tex ends mid-sentence at the bottom of page 76; the
% opening sentence below completes it (top of page 77), still within
% section 2.5.1. Section 2.6 "LU decomposition with row pivoting" begins
% shortly after, on the same page. This file ends mid-equation at the very
% bottom of page 80 (the equation for updating $A(2:5,2:5)$); its
% continuation is on page 81, owned by a different agent -- do not
% fabricate that continuation here.

see (10.76) in Section 10.2.3 for a proof of the last equality.

\section{LU decomposition with row pivoting}

The LU decomposition without pivoting of a matrix has a major drawback: it does not
exist for all nonsingular matrices. In fact, it does not exist even for very well
conditioned matrices. For example, the matrix $A = \begin{pmatrix} 0 & 1\\ 1 & 1
\end{pmatrix}$ does not have an LU decomposition without row pivoting. If it did,
then $A=LU$ could be written as
\begin{equation}
\begin{pmatrix} 1 & 0\\ L(2,1) & 1\end{pmatrix}
\begin{pmatrix} U(1,1) & U(1,2)\\ 0 & U(2,2)\end{pmatrix}
\;=\;
\begin{pmatrix} 0 & 1\\ 1 & 1\end{pmatrix}.
\end{equation}
By multiplying the matrix $U$ by the first row of $L$, we find that $U(1,1)=0$ and
$U(1,2)=1$. Thus, (2.50) becomes
\[
\begin{pmatrix} 1 & 0\\ L(2,1) & 1\end{pmatrix}
\begin{pmatrix} 0 & 1\\ 0 & U(2,2)\end{pmatrix}
\;=\;
\begin{pmatrix} 0 & 1\\ 1 & 1\end{pmatrix}.
\]
However, multiplying the second row of $L$ by the first column of $U$ we obtain that
\[
L(2,1)\cdot 0 + 1\cdot 0 \;=\; 1,
\]
which is not possible. The reason for this is that $A(1,1)=0$, and, since $A(1,1)$
is the first leading principal minor of $A$, the matrix $A$ does not satisfy the
requirement for the existence of the LU decomposition without pivoting from
Theorem 2.1, i.e., that all the leading principal minors are nonzero.

However, any nonsingular matrix has an LU decomposition with row pivoting. To
introduce the concept of pivoting, recall from section 10.1.2 that a permutation
matrix is a matrix obtained by permuting the rows of the identity matrix, and that
matrix multiplication by a permutation matrix to the left results in a
corresponding permutation of the rows of the matrix. For example, the permutation
matrix
\[
P \;=\;
\begin{pmatrix} 0&0&1&0\\ 1&0&0&0\\ 0&0&0&1\\ 0&1&0&0\end{pmatrix}
\quad \text{has row form} \quad
P \;=\; \begin{pmatrix} e_3^t\\ e_1^t\\ e_4^t\\ e_2^t\end{pmatrix},
\]
and is also denoted, using column notation, as
\[
P \;=\; \begin{pmatrix} 3\\ 1\\ 4\\ 2\end{pmatrix}.
\]
If $A$ is a $4\times 4$ matrix with row form
\[
A \;=\; \begin{pmatrix} r_1\\ r_2\\ r_3\\ r_4\end{pmatrix},
\quad \text{then} \quad
PA \;=\; \begin{pmatrix} r_3\\ r_1\\ r_4\\ r_2\end{pmatrix}.
\]

\begin{definition}
The LU decomposition with row pivoting of a nonsingular square matrix $A$ consists
of finding a permutation matrix $P$, a lower triangular matrix $L$ with all entries
on the main diagonal equal to $1$, and a nonsingular upper triangular matrix $U$
such that
\[
PA \;=\; LU.
\]
\end{definition}

\begin{theorem}
Any nonsingular matrix has an LU decomposition with row pivoting.
\end{theorem}

A proof of this result can be found, e.g., in Datta~[12].

For a given permutation matrix $P$, the $L$ and $U$ factors from $PA=LU$ are
unique. However, for any nonsingular matrix $A$, there may be many different
permutation matrices $P$ such that the matrix $PA$ has an LU decomposition. The
algorithm described below, and whose pseudocode is included in Table 2.10,
uniquely identifies the permutation matrix $P$ and the $L$ and $U$ factors for the
LU decomposition with row pivoting of a given nonsingular matrix $A$. Note that all
the entries of the matrix $L$ identified using this algorithm are less than or
equal in absolute value to $1$.

Let $A$ be a nonsingular matrix. We identify a permutation matrix $P$ and a lower
triangular matrix $L$ (with main diagonal entries equal to $1$) and an upper
triangular matrix $U$ as detailed below.

\begin{itemize}
\item Identify the largest entry (in absolute value) in the first column vector
$A(1:n,1)$ of $A$. Switch the corresponding row and the first row of $A$, and do
the same for the permutation matrix $P$, which was initialized to be equal to the
identity matrix:
\end{itemize}

\begin{center}
\fbox{\begin{minipage}{0.85\textwidth}
\raggedright
$P = 1:n$\\
find $i\_max$, the location of the largest entry from $A(1:n,1)$\\
\hspace*{1.5em}$vv = A(1,1:n);\ A(1,1:n) = A(i\_max,1:n);\ A(i\_max,1:n) = vv;$\\
\hspace*{1.5em}$cc = P(1);\ P(1) = P(i\_max);\ P(i\_max) = cc;$\\
clear $vv$ $cc$
\end{minipage}}
\end{center}

\begin{itemize}
\item Compute the first row of $U$ and the first column of $L$:
\end{itemize}

\begin{center}
\fbox{\begin{minipage}{0.5\textwidth}
\raggedright
for $k = 1 : n$\\
\hspace*{1.5em}$U(1,k) \;=\; A(1,k)$\\
\hspace*{1.5em}$L(k,1) \;=\; \dfrac{A(k,1)}{U(1,1)}$\\
end
\end{minipage}}
\end{center}

\begin{itemize}
\item Update the $(n-1)\times(n-1)$ lower right part of $A$ as follows:
\end{itemize}
\begin{equation}
A(2:n,2:n) \;=\; A(2:n,2:n) \;-\; L(2:n,1)\,U(1,2:n),
\end{equation}
which can be written entry by entry as follows:

\begin{center}
\fbox{\begin{minipage}{0.5\textwidth}
\raggedright
for $j = 2 : n$\\
\hspace*{1.5em}for $k = 2 : n$\\
\hspace*{3em}$A(j,k) = A(j,k) - L(j,1)U(1,k)$\\
\hspace*{1.5em}end\\
end
\end{minipage}}
\end{center}

The row permutation process as well as the computation of every row of $U$ and
column of $L$ are done recursively from the latest updated part of the matrix $A$;
see also the $5\times 5$ example below. For example, to compute the $i$-th row of
$U$ and the $i$-th column of $L$, we do the following:

\begin{itemize}
\item Identify the largest entry (in absolute value) in the first column vector
$A(i:n,i)$ of $A(i:n,i:n)$. Switch the corresponding row and the first row of
$A(i:n,i:n)$. Switch the corresponding row and the row $i$ of the permutation
matrix $P$. Also, switch all the entries already computed from the corresponding
row and the row $i$ of the matrix $L$:
\end{itemize}

\begin{center}
\fbox{\begin{minipage}{0.85\textwidth}
\raggedright
find $i\_max$, the location of the largest entry from $A(i:n,i)$\\
$vv = A(i,i:n);\ A(i,i:n) = A(i\_max,i:n);\ A(i\_max,i:n) = vv;$\\
$cc = P(i);\ P(i) = P(i\_max);\ P(i\_max) = cc;$\\
\hspace*{1.5em}$ww = L(i,1:(i-1));$\\
\hspace*{1.5em}$L(i,1:(i-1)) = L(i\_max,1:(i-1));\ L(i\_max,1:(i-1)) = ww;$\\
clear $vv$ $ww$ $cc$
\end{minipage}}
\end{center}

\begin{itemize}
\item Compute the $i$-th row of $U$ and the $i$-th column of $L$:
\end{itemize}

\begin{center}
\fbox{\begin{minipage}{0.5\textwidth}
\raggedright
for $k = i : n$\\
\hspace*{1.5em}$U(i,k) \;=\; A(i,k)$\\
\hspace*{1.5em}$L(k,i) \;=\; \dfrac{A(k,i)}{U(i,i)}$\\
end
\end{minipage}}
\end{center}

\begin{itemize}
\item Update the $(n-i)\times(n-i)$ lower right part of $A$ as follows:
\end{itemize}
\begin{equation}
A(i+1:n,i+1:n) \;=\; A(i+1:n,i+1:n) \;-\; L(i+1:n,i)\,U(i,i+1:n),
\end{equation}
which can be written entry by entry as follows:

\begin{center}
\fbox{\begin{minipage}{0.55\textwidth}
\raggedright
for $j = (i+1) : n$\\
\hspace*{1.5em}for $k = (i+1) : n$\\
\hspace*{3em}$A(j,k) = A(j,k) - L(j,i)U(i,k)$\\
\hspace*{1.5em}end\\
end
\end{minipage}}
\end{center}

Further clarification on the recursive part of the LU decomposition with row
pivoting can be found in the example below for a $5\times 5$ matrix.

% NOTE: this worked example is not finished within pages 74-80 -- it
% continues past the bottom of page 80 onto page 81, which belongs to a
% different agent's file. The \examplex environment (see preamble.tex) is
% therefore deliberately NOT used here, since closing it would print the
% closing square before the example actually ends. The continuation file
% should open with \textit{Example:} already implied (do not print a
% second "Example:" label) and must supply the closing
% \hfill$\square$\par\medskip once the example concludes.
\par\medskip\noindent\textit{Example:}\ Let
\[
A \;=\; \begin{pmatrix}
1 & 2 & -7 & -1.5 & 2\\
4 & 4 & 0 & -6 & -2\\
-2 & -1 & 2 & 6 & 2.5\\
0 & -2 & 2 & -4 & 1\\
2 & 1 & 13 & -5 & 3.5
\end{pmatrix}.
\]
Let $L$ and $U$ be the LU factors from the LU decomposition with row pivoting of
$A$ corresponding to a permutation matrix $P$ which will be identified below, i.e.,
such that $PA=LU$.

The permutation matrix $P$ is set equal to the identity matrix and will be updated
as rows are switched. We will use the following notation for the row form of the
permutation matrix $P$:
\[
P \;=\; \begin{pmatrix} 1\\ 2\\ 3\\ 4\\ 5\end{pmatrix}
\;=\;
\begin{pmatrix} e_1^t\\ e_2^t\\ e_3^t\\ e_4^t\\ e_5^t\end{pmatrix}.
\]

We begin by identifying the largest entry, in absolute value, from the first
column
\[
\begin{pmatrix} 1\\ 4\\ -2\\ 0\\ 2\end{pmatrix}
\]
of $A$. That entry is $4$, corresponding to the second row of $A$. We switch the
first row and the second row of $A$. The new form of the matrix $A$ is
\begin{equation}
A \;=\; \begin{pmatrix}
4 & 4 & 0 & -6 & -2\\
1 & 2 & -7 & -1.5 & 2\\
-2 & -1 & 2 & 6 & 2.5\\
0 & -2 & 2 & -4 & 1\\
2 & 1 & 13 & -5 & 3.5
\end{pmatrix}.
\end{equation}

We record the switch by permuting the first row and the second row of $P$ as
follows:
\begin{equation}
P \;=\; \begin{pmatrix} 1\\ 2\\ 3\\ 4\\ 5\end{pmatrix}
\quad \text{becomes} \quad
P \;=\; \begin{pmatrix} \mathbf{2}\\ \mathbf{1}\\ 3\\ 4\\ 5\end{pmatrix}
\;=\;
\begin{pmatrix} e_2^t\\ e_1^t\\ e_3^t\\ e_4^t\\ e_5^t\end{pmatrix}.
\end{equation}

After the row permutations are done, the first row of $U$ and the first column of
$L$ are computed as in the algorithm for LU decomposition without row pivoting.

The entries of first row of $U$ are given by (2.22), i.e., $U(1,k) = A(1,k)$ for
$k=1:5$, where the matrix $A$ is given by (2.53):
\[
U(1,1) = 4;\quad U(1,2) = 4;\quad U(1,3) = 0;
\]
\[
U(1,4) = -6;\quad U(1,5) = -2.
\]
The entries of first column of $L$ are given by (2.24), i.e., $L(k,1) =
\dfrac{A(k,1)}{U(1,1)}$ for $k=1:5$, where the matrix $A$ is given by (2.53):
\[
L(1,1) = 1;\quad L(2,1) = \frac{1}{U(1,1)} = \frac{1}{4} = 0.25;
\]
\[
L(3,1) = \frac{-2}{U(1,1)} = \frac{-2}{4} = -0.5;
\]
\[
L(4,1) = \frac{0}{U(1,1)} = 0;\quad L(5,1) = \frac{2}{U(1,1)} = \frac{2}{4} = 0.5.
\]

The current forms of $L$ and $U$ are
\begin{equation}
L \;=\; \begin{pmatrix} 1\\ 0.25\\ -0.5\\ 0\\ 0.5\end{pmatrix};
\qquad
U \;=\; \begin{pmatrix}
4 & 4 & 0 & -6 & -2\\
0 & U(2,2) & U(2,3) & U(2,4) & U(2,5)\\
0 & 0 & U(3,3) & U(3,4) & U(3,5)\\
0 & 0 & 0 & U(4,4) & U(4,5)\\
0 & 0 & 0 & U(5,4) & U(5,5)
\end{pmatrix}.
\end{equation}
Note that, for the current form of the matrix $L$, we only consider the entries of
$L$ computed so far, since the rows of $L$ will switch as the corresponding rows of
$A$ switch.

The updated form of the $4\times 4$ matrix $A(2:5,2:5)$ is computed using (2.51),
from $A(2:5,2:5)$ obtained from (2.53) and with $L$ and $U$ given by (2.55), as
follows:
\begin{align*}
& A(2:5,2:5)\\
& \;=\; A(2:5,2:5) \;-\; L(2:5,1)\, U(1,2:5)
\end{align*}
% Example continues on page 81 (owned by a different agent); the closing
% \hfill$\square$\par\medskip belongs in that continuation file, not here.
